%
% tabular.tex -- table using xcolor
%
% (c) 2021 Prof Dr Andreas Müller, OST Ostschweizer Fachhochschule
%
\documentclass[tikz]{standalone}
\usepackage{amsmath}
\usepackage{times}
\usepackage{txfonts}
\usepackage{pgfplots}
\usepackage{csvsimple}
\usepackage[table]{xcolor}
\usetikzlibrary{arrows,intersections,math,calc}
\definecolor{darkred}{rgb}{0.8,0,0}
\definecolor{darkgreen}{rgb}{0,0.6,0}
\begin{document}
\def\skala{1}
\begin{tikzpicture}[>=latex,thick,scale=\skala]

\node at (0,0) {\begin{tabular}{|c|c|c|}
	\hline
	\rowcolor{blue!40}
	\textcolor{white}{\textbf{\textit{s}}} & 
	\textcolor{white}{\textbf{$\zeta(s)$}} & 
	\textcolor{white}{\textbf{Ungefährer Wert}\raisebox{2pt}{\strut}} \\
	\hline
	$2$ & $\pi^2/6$ & $\approx 1{,}6449$ \raisebox{2pt}{\strut}\\
%	\hline
	$4$ & $\pi^4/90$ & $\approx 1{,}0823$ \raisebox{2pt}{\strut}\\
%	\hline
	$0$ & $-1/2$ & (analytische Fortsetzung) \raisebox{2pt}{\strut}\\
	\hline
\end{tabular}};

\end{tikzpicture}
\end{document}

